% =========================
% Archivo: introduccion.tex
% =========================

\section{Introducción}

La inteligencia artificial permite diseñar sistemas capaces de tomar decisiones, optimizar recursos y actuar de manera autónoma en entornos complejos. En áreas como la logística y la automatización, estos sistemas resultan especialmente importantes debido a la necesidad de coordinar rutas, administrar recursos y operar robots en situaciones dinámicas e inciertas.

En el presente trabajo se analiza el caso de una empresa de logística y automatización que desea implementar un sistema inteligente para resolver distintos problemas relacionados con distribución y operación de almacenes. El objetivo principal es identificar qué algoritmos o enfoques de inteligencia artificial resultan más adecuados para cada situación planteada.

Para ello, el problema se divide en cuatro partes principales. Primero, se estudian problemas de optimización y búsqueda local aplicados a rutas entre ciudades. Luego, se aborda la búsqueda local en espacios continuos mediante la ubicación de puntos de servicio o recarga en un mapa. Posteriormente, se analiza la búsqueda con acciones no deterministas en el contexto de un robot cuyas acciones pueden producir diferentes resultados. Finalmente, se trabaja el razonamiento en entornos parcialmente observables, donde el robot opera con información limitada obtenida por sensores.

A lo largo del informe se realizará el modelado básico de cada problema, identificando estados, acciones, objetivos y dificultades principales. Además, se justificará la elección de cada algoritmo o estrategia, comparando sus ventajas y limitaciones según el tipo de problema.

\subsection{Objetivo general}

Analizar y comparar distintos enfoques de búsqueda y optimización en inteligencia artificial, aplicándolos a un entorno de logística y automatización para determinar qué estrategia resulta más adecuada en cada caso.

\subsection{Objetivos específicos}

\begin{itemize}
    \item Modelar problemas de optimización utilizando búsqueda local.
    
    \item Analizar situaciones con variables continuas y explicar cómo resolverlas mediante técnicas de optimización continua.
    
    \item Representar problemas con acciones no deterministas mediante planes condicionales y árboles AND-OR.
    
    \item Explicar el razonamiento de un agente en entornos parcialmente observables utilizando estados de creencia y percepciones.
    
    \item Comparar las estrategias utilizadas y justificar su aplicación en cada situación.
\end{itemize}

\subsection{Estructura del informe}

El informe se encuentra organizado en cuatro bloques principales, cada uno correspondiente a un tipo de problema planteado en el caso de estudio. En cada sección se describe el problema, se realiza su modelado básico y se analiza el algoritmo o enfoque más adecuado para resolverlo. Finalmente, se presenta una comparación general entre las estrategias utilizadas y las conclusiones obtenidas.